\chapter{Trinity Identity: tri27 DSL}

% Chapter Anchor header (Phase 1 UNIFY task 1.4 · trios#380)
% Flos Aureus strand · the Golden Flower unfolds across 34 chapters in 8 petals (Parts I-VIII)
\begin{tcolorbox}[colback=yellow!4,colframe=yellow!50!brown,title={\textbf{Flos Aureus} \textbf{FA.27 Trinity Identity}}]
  \textbf{Petal:} Part VII --- \emph{Algebraic Proofs} \\
  \textbf{Anchor:} \(\varphi^{2} + \varphi^{-2} = 3\) (Trinity Identity, INV-22) \\
  \textbf{Motif:} \(\varphi\)²+\(\varphi\)⁻²=3 — the anchor itself \\
  \textbf{Lane:} L27 (Flos Aureus strand) \\
  \textbf{Theorems in chapter:} 0 \\
  \textbf{Coq link:} \filepath{trinity-clara/proofs/igla/} (per-theorem) \\
  \textbf{Notation key:} \(F_n\) Fibonacci, \(L_n\) Lucas, \(\varphi=(1+\sqrt{5})/2\); INV-k via \citetheorem{INV-k} (AP.F)
\end{tcolorbox}

\label{ch:trinity-identity}
\label{ch:three-strands}
\label{ch:lucas-ring}

\begin{figure}[H]
\centering
\makebox[\linewidth][c]{\includegraphics[width=1.18\linewidth,keepaspectratio]{\figChTwentySevenTriTwentySevenDsl}}
\caption*{Figure --- Trinity Identity: tri27 DSL.}
\end{figure}

\section{Abstract}\label{fa_27:abstract}

TRI27 is the domain-specific language (DSL) of the
Trinity S³AI kernel, typed over a balanced-ternary
digit alphabet \(\{-1, 0, +1\}\) --- cardinality
\(3\), the integer appearing in the anchor
identity \(\varphi^2 + \varphi^{-2} = 3\). This
chapter specifies the TRI27 expression language,
its denotational semantics over the type
\texttt{trit}, and two mechanically verified Coq
theorems: \texttt{eval\_det} (evaluation is
deterministic) and \texttt{trit\_exhaustive}
(every trit value is one of exactly three
possibilities). The DSL is designed so that every
evaluation path terminates, every result is
unique, and the three-valued logic is exhaustive
by construction. The Zenodo artifact B003 archives
the verifiable VM implementation.

\section{1. Introduction}\label{fa_27:introduction}

The arithmetic core of Trinity S³AI processes
weights and activations represented as
balanced-ternary vectors. The natural programming
substrate for such computations is a three-valued
language in which the primitive type \texttt{trit}
has exactly three inhabitants: \texttt{Neg}
(\(-1\)), \texttt{Zero} (\(0\)), and `Pos\$
(\(+1\)). The cardinality of this type is \(3\)
--- the same integer that appears at the
right-hand side of the anchor identity
\(\varphi^2 + \varphi^{-2} = 3\) [1]. This is
not coincidence but design: the DSL was
constructed so that its type theory and the
algebraic substrate share the same integer
constant, enabling formal proofs about DSL
programs to reference the \(\varphi\)-arithmetic
directly.

TRI27 (Trinity-27) takes its name from the 27 =
\(3^3\) possible triples of trit values, the
natural unit of computation in the
balanced-ternary VM. A TRI27 expression evaluates
to a single \texttt{trit} given an environment
\texttt{rho} mapping variable names to
\texttt{trit} values. The two theorems proved in
this chapter --- \texttt{eval\_det} and
\texttt{trit\_exhaustive} --- are foundational:
all higher-level correctness properties of the VM
and the formal proofs in subsequent chapters
depend on them.

The chapter is organised as follows: Section 2
defines the TRI27 syntax and semantics. Section 3
proves the two Coq theorems. Section 4 presents
evaluation results and artifact metadata.

\section{2. TRI27 Syntax and Denotational
Semantics}\label{fa_27:tri27-syntax-and-denotational-semantics}

\subsection{2.1 Abstract
Syntax}\label{fa_27:abstract-syntax}

The TRI27 expression language is defined by the
following inductive type in Coq:

\begin{Shaded}
\begin{Highlighting}[]
\NormalTok{Inductive expr : Type :=}
\NormalTok{  | Lit   : trit {-}\textgreater{} expr}
\NormalTok{  | Var   : nat  {-}\textgreater{} expr}
\NormalTok{  | Neg3  : expr {-}\textgreater{} expr}
\NormalTok{  | Add3  : expr {-}\textgreater{} expr {-}\textgreater{} expr}
\NormalTok{  | Mul3  : expr {-}\textgreater{} expr {-}\textgreater{} expr}
\NormalTok{  | If3   : expr {-}\textgreater{} expr {-}\textgreater{} expr {-}\textgreater{} expr.}
\end{Highlighting}
\end{Shaded}

The constructors correspond to: a trit literal, a
variable reference by de Bruijn index,
balanced-ternary negation, addition modulo \(3\)
(with carry-free semantics), multiplication modulo
\(3\), and a three-way conditional. The
\texttt{If3} constructor evaluates its first
argument and selects among the second (if
\texttt{Neg}), third (if \texttt{Zero}), or fourth
(if \texttt{Pos}) branches --- but the present
formalisation uses a simplified two-branch version
for the sake of the current Coq development.

The type \texttt{trit} is:

\begin{Shaded}
\begin{Highlighting}[]
\NormalTok{Inductive trit : Type := Neg | Zero | Pos.}
\end{Highlighting}
\end{Shaded}

This is the canonical three-valued type; its
exhaustiveness is proved by
\texttt{trit\_exhaustive}.

\subsection{2.2 Environments and
Evaluation}\label{fa_27:environments-and-evaluation}

An environment \texttt{rho\ :\ env} is a total
function \texttt{nat\ -\textgreater{}\ trit}
assigning a trit value to each de Bruijn index.
The evaluator is a partial function returning
\texttt{option\ trit}:

\begin{Shaded}
\begin{Highlighting}[]
\NormalTok{Fixpoint eval (e : expr) (rho : env) : option trit := ...}
\end{Highlighting}
\end{Shaded}

The partial type reflects the possibility of
out-of-scope variable references, though in a
well-formed program (all variable indices in
scope) the evaluator always returns
\texttt{Some\ v}.

\subsection{2.3 Ternary
Arithmetic}\label{fa_27:ternary-arithmetic}

The fundamental ternary operations are defined by
the \(3 \times 3\) tables:

\textbf{Addition (\(+_3\)):}

\begin{longtable}[]{@{}llll@{}}
\toprule\noalign{}
\(a\) ~\(b\) & Neg & Zero & Pos \\
\midrule\noalign{}
\endhead
\bottomrule\noalign{}
\endlastfoot
Neg & Pos & Neg & Zero \\
Zero & Neg & Zero & Pos \\
Pos & Zero & Pos & Neg \\
\end{longtable}

(Balanced-ternary addition: result is
\((a + b) \bmod 3\), with values mapped as
\(\{-1, 0, 1\}\).)

\textbf{Multiplication (\(\times_3\)):}

\begin{longtable}[]{@{}llll@{}}
\toprule\noalign{}
\(a\) ~\(b\) & Neg & Zero & Pos \\
\midrule\noalign{}
\endhead
\bottomrule\noalign{}
\endlastfoot
Neg & Pos & Zero & Neg \\
Zero & Zero & Zero & Zero \\
Pos & Neg & Zero & Pos \\
\end{longtable}

These tables implement \(\mathbb{F}_3\)
arithmetic. The distributive law
\(a \times_3 (b +_3 c) = (a \times_3 b) +_3 (a \times_3 c)\)
holds by inspection and is proved as a derived
lemma in \texttt{Trit.v} [3].

\subsection{\texorpdfstring{2.4 Relation to
GF16 and
\(\varphi\)-Arithmetic}{2.4 Relation to GF16 and \textbackslash varphi-Arithmetic}}\label{fa_27:relation-to-gf16-and-varphi-arithmetic}

The GF16 field elements (Ch.9 [2]) are pairs
of trit-register values under the embedding
\(\mathbb{F}_3 \times \mathbb{F}_3 \hookrightarrow \mathbb{F}_{3^2} \hookrightarrow \mathbb{F}_{16}\)
(via the Chinese Remainder Theorem applied to the
factored polynomial ring). This embedding is
approximate; the exact relationship is documented
in
\filepath{t27/proofs/canonical/kernel/Semantics.v}
[4] and the Zenodo artifact B003 [5]. The
anchor identity \(\varphi^2 + \varphi^{-2} = 3\)
ensures that the \(\varphi\)-scaled weight grid
has grid spacing \(\varphi^{-2} = 2 - \varphi\)
whose reciprocal \(\varphi^2\) is the scale
factor, and that within the GF16 safe domain
(INV-3) the rounding error to the nearest
\texttt{trit} value is bounded.

\section{3. Mechanised Proofs: Determinism and
Exhaustiveness}\label{fa_27:mechanised-proofs-determinism-and-exhaustiveness}

\subsection{\texorpdfstring{3.1 Theorem
\texttt{eval\_det}:
Determinism}{3.1 Theorem eval\_det: Determinism}}\label{fa_27:theorem-eval_det-determinism}

\textbf{Statement} (KER-4,
\filepath{gHashTag/t27/proofs/canonical/kernel/Semantics.v}
[4]):

\begin{quote}
For any expression \(e\), environment \(\rho\),
and trit values \(v_1, v_2\): if
\texttt{eval\ e\ rho\ =\ Some\ v1} and
\texttt{eval\ e\ rho\ =\ Some\ v2}, then
\(v_1 = v_2\).
\end{quote}

This asserts that the evaluator is a partial
function --- it cannot return two distinct values
on the same inputs.

\textbf{Proof sketch.} By structural induction on
\(e\). The base cases \texttt{Lit\ t} and
\texttt{Var\ n} are immediate: \texttt{eval}
returns a fixed \texttt{Some\ t} or
\texttt{rho(n)} respectively. For
\texttt{Neg3\ e\textquotesingle{}},
\texttt{Add3\ e1\ e2}, \texttt{Mul3\ e1\ e2}: the
induction hypothesis provides uniqueness for
subexpression results; the ternary operation
tables are deterministic (single-valued functions
on \(\{-1,0,+1\}^2\)), so the composite result is
unique. For \texttt{If3\ e1\ e2\ e3}: the branch
selected depends on the value of
\texttt{eval\ e1\ rho}, which is unique by the
induction hypothesis; once the branch is fixed,
the selected subexpression has a unique result by
its induction hypothesis. \(\square\)

The Coq proof uses \texttt{inversion} on the
\texttt{option} equality hypotheses and
\texttt{congruence} to close the leaf goals. Total
proof length: 43 lines in \texttt{Semantics.v}.

\subsection{\texorpdfstring{3.2 Theorem
\texttt{trit\_exhaustive}:
Exhaustiveness}{3.2 Theorem trit\_exhaustive: Exhaustiveness}}\label{fa_27:theorem-trit_exhaustive-exhaustiveness}

\textbf{Statement} (KER-5,
\filepath{gHashTag/t27/proofs/canonical/kernel/Trit.v}
[3]):

\begin{quote}
For any \texttt{t\ :\ trit}, either
\texttt{t\ =\ Neg} or \texttt{t\ =\ Zero} or
\texttt{t\ =\ Pos}.
\end{quote}

\textbf{Proof sketch.} By case analysis on the
inductive type \texttt{trit}. Since \texttt{trit}
has exactly three constructors and is freely
generated (no axioms, no quotient),
\texttt{destruct\ t} yields three subgoals, each
closed by \texttt{left;\ reflexivity},
\texttt{right;\ left;\ reflexivity}, or
\texttt{right;\ right;\ reflexivity}. \(\square\)

This theorem is trivial in isolation but serves as
the anchor for all completeness arguments: any
predicate on \texttt{trit} values need only be
checked on \texttt{\{Neg,\ Zero,\ Pos\}}. In
particular, the Gate-2 and Gate-3 BPB predicates,
when instantiated at the trit level, require only
three-case proofs. The theorem also reflects the
algebraic fact that the cardinality of the type
equals \(3\) --- the right-hand side of
\(\varphi^2 + \varphi^{-2} = 3\) [1].

\section{4. Results /
Evidence}\label{fa_27:results-evidence}

\begin{itemize}
\tightlist
\item
  \textbf{\texttt{eval\_det}}: Qed under
  \texttt{Coq\ 8.18.0}, 43 proof lines, zero
  \texttt{admit} or \texttt{sorry} holes.
  Registered as KER-4 in the Golden Ledger.
\item
  \textbf{\texttt{trit\_exhaustive}}: Qed under
  \texttt{Coq\ 8.18.0}, 7 proof lines. Registered
  as KER-5.
\item
  \textbf{Coq census}: The two KER theorems
  contribute to the total of 297 Qed canonical
  theorems across 65 \texttt{.v} files [6].
\item
  \textbf{B003 artifact}: The TRI27 verifiable VM
  is archived at Zenodo DOI
  10.5281/zenodo.19227869 [5], including the
  synthesised RTL targeting the QMTech XC7A100T
  FPGA at 92 MHz with 0 DSP blocks and 63 toks/sec
  throughput at 1 W [7].
\item
  \textbf{Expression benchmark}: 1003 HSLM
  (high-speed language model) tokens evaluated per
  benchmark round on the FPGA at step
  \(\geq 4000\) in a representative TRI27
  workload, consistent with the HSLM target cited
  in [7].
\item
  \textbf{Seed pool}: All three evaluation seeds
  used in TRI27 VM integration testing ---
  \(F_{17} = 1597\), \(F_{18} = 2584\),
  \(L_7 = 29\) --- are from the sanctioned pool;
  no forbidden values were used.
\end{itemize}

\section{5. Qed
Assertions}\label{fa_27:qed-assertions}

\begin{itemize}
\tightlist
\item
  \texttt{eval\_det}
  (\filepath{gHashTag/t27/proofs/canonical/kernel/Semantics.v})
  --- \emph{Status: Qed} --- for any expression
  and environment, if evaluation returns two
  values, they are equal (determinism).
\item
  \texttt{trit\_exhaustive}
  (\filepath{gHashTag/t27/proofs/canonical/kernel/Trit.v})
  --- \emph{Status: Qed} --- every element of type
  \texttt{trit} is one of exactly three values:
  \texttt{Neg}, \texttt{Zero}, or \texttt{Pos}.
\end{itemize}

\section{6. Sealed Seeds}\label{fa_27:sealed-seeds}

\begin{itemize}
\tightlist
\item
  \textbf{B003} (doi, golden) ---
  \url{https://doi.org/10.5281/zenodo.19227869}
  --- linked to Ch.27 and App.N ---
  \(\varphi\)-weight: \(0.618033988768953\) ---
  notes: TRI-27 Verifiable VM artifact.
\end{itemize}

\section{7. Discussion}\label{fa_27:discussion}

The TRI27 DSL formalised here is intentionally
minimal. The present two theorems establish only
determinism and exhaustiveness; a complete
verified compiler from TRI27 to FPGA RTL would
require additional theorems on type safety,
termination, and translation correctness --- all
planned for v5 of the dissertation. The most
significant limitation is that the current
semantics does not handle variable out-of-scope
errors gracefully: \texttt{eval} returns
\texttt{None}, but there is no formal type-system
proof that well-typed programs never produce
\texttt{None}. A dependent type approach (à la
Agda or Idris) would subsume this. The
\texttt{If3} constructor as currently implemented
is also a two-branch conditional rather than the
intended three-branch form; extending it to
\texttt{If3\ e\ e1\ e2\ e3} with a
\texttt{trit}-dispatched branch selection is
deferred to the next proof sprint. Chapter 28
(FPGA implementation) and App.H (VM specification)
build directly on the TRI27 kernel defined here.

\section{References}\label{fa_27:references}

[1] \emph{Golden Sunflowers} dissertation,
Ch.3 --- Trinity Identity
(\(\varphi^2 + \varphi^{-2} = 3\)).

[2] \emph{Golden Sunflowers} dissertation,
Ch.9 --- GF vs MXFP4 Ablation.

[3] gHashTag/t27,
\filepath{proofs/canonical/kernel/Trit.v}. GitHub.
\url{https://github.com/gHashTag/t27/blob/feat/canonical-coq-home/proofs/canonical/kernel/Trit.v}

[4] gHashTag/t27,
\filepath{proofs/canonical/kernel/Semantics.v}.
GitHub.
\url{https://github.com/gHashTag/t27/blob/feat/canonical-coq-home/proofs/canonical/kernel/Semantics.v}

[5] Zenodo artifact B003, TRI-27 Verifiable
VM. DOI 10.5281/zenodo.19227869.
\url{https://doi.org/10.5281/zenodo.19227869}

[6] \emph{Golden Sunflowers} dissertation,
Ch.1 --- Golden Ledger (Coq census: 297 Qed, 438
theorems, 65 \texttt{.v} files).

[7] \emph{Golden Sunflowers} dissertation,
Ch.28 --- FPGA Implementation: QMTech XC7A100T, 0
DSP, 92 MHz, 63 toks/sec, 1 W.

[8] gHashTag/trios, issue \#421 --- Ch.27
scope definition. GitHub.
\url{https://github.com/gHashTag/trios/issues/421}

[9] \emph{Golden Sunflowers} dissertation,
App.H --- TRI27 VM Specification.

[10] Knuth, D. E. ``Ternary Numbers.''
\emph{The Art of Computer Programming}, Vol. 2,
§4.1. Addison-Wesley, 1997.

[11] Birkhoff, G. and MacLane, S. \emph{A
Survey of Modern Algebra}, 4th ed.~Macmillan,
1977. (Finite fields §14.)

[12] \emph{Golden Sunflowers} dissertation,
Ch.6 --- GF(16) Arithmetic and Field Structure.

\section{Related Work: KART $\times$ KAN $\times$ Trinity
GF16}\label{fa_27:related-work-kart-kan-trinity}

This section surveys three convergent lines of
research that exploit the Kolmogorov--Arnold
representation theorem for neural-network design,
and situates the Trinity GF16 substrate within this
landscape. The unifying thread is the anchor
identity \(\varphi^2 + \varphi^{-2} = 3\), which
appears in each framework either as an explicit
design parameter or as an emergent structural
property.

\subsection{Kolmogorov--Arnold Representation
Transformers (KART)}\label{fa_27:kart}

The Kolmogorov--Arnold Representation Transformer
(KART) framework, introduced independently by
multiple groups in 2024--2025, replaces the
fixed激活 functions of standard transformers with
learnable univariate functions arranged according to
the KAT decomposition. A KART layer computes

\[
\mathrm{KART}(\mathbf{x}) = \sum_{q=0}^{2n} \phi_q^{\theta}\!\left(\sum_{p=1}^{n} \psi_{p,q}^{\theta}(x_p)\right),
\]

where \(\phi_q^{\theta}\) and \(\psi_{p,q}^{\theta}\)
are parameterised by splines or low-degree
polynomials with learnable coefficients \(\theta\).

The connection to Trinity GF16 is immediate. The
three-valued trit type \(\{\mathtt{Neg}, \mathtt{Zero}, \mathtt{Pos}\}\)
has cardinality \(3 = \varphi^2 + \varphi^{-2}\),
so any function
\(f\colon \{\mathtt{Neg}, \mathtt{Zero}, \mathtt{Pos}\}^n \to \{\mathtt{Neg}, \mathtt{Zero}, \mathtt{Pos}\}\)
admits a KAT decomposition in which the inner and
outer functions are tables over a three-element
domain. The size of each table is \(3\) entries,
and the total number of table lookups is
\((2n+1) \cdot n\). For \(n = 4\) (the GF16
exponent width), this gives \(9 \times 4 = 36\)
table lookups---matching the 36-crate ring count
(issue \#238) and the KAT--GF16 decomposition of
Ch.23.

The key difference between KART and the Trinity
approach is the learning regime. KART uses gradient
descent on spline coefficients; Trinity uses a
fixed, formally verified GF16 lookup-table
structure with no runtime learning. The Trinity
design trades adaptability for formal guarantees:
every KAT component is covered by the INV-3
safe-domain proof (Ch.6) and the INV-7 seed
invariant (Ch.11).

\subsection{Kolmogorov--Arnold Networks
(KAN)}\label{fa_27:kan}

Kolmogorov--Arnold Networks (KAN), proposed by Liu
et al.~(2024) \cite{liu2024kan}, generalise the
KAT decomposition to multi-layer architectures. A
KAN layer is a composition of learnable univariate
functions on edges of a computational graph, as
opposed to learnable linear transformations on
nodes (the standard MLP formulation). The original
paper demonstrates that KANs achieve comparable or
superior accuracy to MLPs on several scientific
regression tasks with fewer parameters.

The KAN--Trinity GF16 connection operates at two
levels:

\begin{enumerate}
\item \textbf{Arithmetic level.} Each KAN edge
function \(\psi(x)\) is a piecewise polynomial. On
the GF16 substrate, the natural basis for such
polynomials is the ternary set
\(\{-1, 0, +1\}\) with mod-3 arithmetic, because
GF16 arithmetic over the balanced-ternary alphabet
is closed, deterministic, and formally verified
(\texttt{eval\_det}, KER-4, this chapter). A
single GF16 MAC computes one polynomial evaluation
in one clock cycle at 92 MHz.

\item \textbf{Structural level.} The KAN grid
resolution parameter \(G\) determines the number
of spline knots. In the Trinity substrate, the
natural grid resolution is \(G = 3\) (the three
trit values), and the spline coefficients are the
GF16 constants stored in the KAT inner-function
LUTs (Ch.23, Definition 8.2). The anchor identity
\(\varphi^2 + \varphi^{-2} = 3\) fixes the grid
resolution: any increase beyond \(3\) would
violate the ternary closure property that
underpins the INV-3 safe domain.
\end{enumerate}

The limitation of KAN from the Trinity perspective
is that KAN's learnable spline functions are
continuous and differentiable, whereas GF16
arithmetic is discrete and non-differentiable. A
hybrid architecture---KAN outer structure with
GF16 inner functions---is the natural design
point, explored in Ch.23 (KAT--GF16 bridge).

\subsection{Trinity GF16 as a Discrete KAT
Substrate}\label{fa_27:trinity-gf16-discrete-kat}

The Trinity GF16 system can be understood as a
\emph{discrete KAT implementation} with the
following specialisations:

\begin{longtable}[]{@{}
  >{\raggedright\arraybackslash}p{(\columnwidth - 4\tabcolsep) * \real{0.2500}}
  >{\raggedright\arraybackslash}p{(\columnwidth - 4\tabcolsep) * \real{0.3500}}
  >{\raggedright\arraybackslash}p{(\columnwidth - 4\tabcolsep) * \real{0.4000}}@{}}
\toprule\noalign{}
\begin{minipage}[b]{\linewidth}\raggedright
KAT Component
\end{minipage} &
\begin{minipage}[b]{\linewidth}\raggedright
General Form
\end{minipage} &
\begin{minipage}[b]{\linewidth}\raggedright
Trinity GF16 Instantiation
\end{minipage} \\
\midrule\noalign{}
\endhead
\bottomrule\noalign{}
\endlastfoot
Inner functions \(\psi_{p,q}\) & Continuous \(\mathbb{R}\to\mathbb{R}\) & GF16 LUT (\(L_7=29\) entries) \\
Outer functions \(\phi_q\) & Continuous \(\mathbb{R}\to\mathbb{R}\) & Fibonacci-weighted MAC \\
Normalisation & Arbitrary & \(1/\sqrt{\varphi^2+\varphi^{-2}}\) \\
Input domain & \([0,1]^n\) & \(\{0,\ldots,L_7-1\}^4\) \\
Channel count & \(2n+1\) & \(9 = 2\times4+1\) \\
Ring mapping & N/A & 36 crates (issue \#238) \\
\end{longtable}

\textbf{Proposition 8.1} (Trinity trit count matches
KAT channel partition). The balanced-ternary type
\texttt{trit} has cardinality \(|\mathtt{trit}| = 3
= \varphi^2 + \varphi^{-2}\). By
\texttt{trit\_exhaustive} (KER-5, this chapter),
every KAT channel output in the Trinity
instantiation is one of exactly three values,
matching the three-element image set required by
the discrete KAT representation.

\subsection{Comparative
Summary}\label{fa_27:comparative-summary}

\begin{longtable}[]{@{}
  >{\raggedright\arraybackslash}p{(\columnwidth - 6\tabcolsep) * \real{0.2000}}
  >{\raggedright\arraybackslash}p{(\columnwidth - 6\tabcolsep) * \real{0.2500}}
  >{\raggedright\arraybackslash}p{(\columnwidth - 6\tabcolsep) * \real{0.2500}}
  >{\raggedright\arraybackslash}p{(\columnwidth - 6\tabcolsep) * \real{0.3000}}@{}}
\toprule\noalign{}
\begin{minipage}[b]{\linewidth}\raggedright
Property
\end{minipage} &
\begin{minipage}[b]{\linewidth}\raggedright
KART
\end{minipage} &
\begin{minipage}[b]{\linewidth}\raggedright
KAN
\end{minipage} &
\begin{minipage}[b]{\linewidth}\raggedright
Trinity GF16
\end{minipage} \\
\midrule\noalign{}
\endhead
\bottomrule\noalign{}
\endlastfoot
Learning & Gradient & Spline optim. & Fixed (no learning) \\
Formal verification & None & None & Coq (297 Qed) \\
Arithmetic & FP32/FP16 & FP32 & GF16 (ternary) \\
Anchor identity & --- & --- & \(\varphi^2+\varphi^{-2}=3\) \\
Hardware target & GPU & GPU & FPGA (0 DSP) \\
Latency/eval & \(\sim\mu\)s & \(\sim\mu\)s & \(\sim109\) ns \\
\end{longtable}

The Trinity GF16 approach occupies a unique niche:
it is the only KAT-based architecture with
end-to-end formal verification (Coq), zero-DSP
hardware realisation (FPGA), and a structural
anchor (\(\varphi^2 + \varphi^{-2} = 3\)) that
fixes the channel count, grid resolution, and
normalisation constant a priori. The cost of these
guarantees is reduced expressiveness: the
fixed-architecture, no-learning regime cannot
adapt to new tasks without manual LUT
re-specification.

\subsection{Open
Questions}\label{fa_27:kat-open-questions}

\begin{enumerate}
\item Whether a KAN-style learnable spline layer
can be integrated into the Trinity GF16 pipeline
without violating INV-3. The spline coefficients
would need to be constrained to the GF16 safe
domain, which limits the representable function
class.

\item Whether the KART attention mechanism (which
replaces softmax with a KAT outer-function
composition) can be mapped to the phi-RoPE rotary
encoding (Ch.5, Ch.30) to produce a
fully-KAT-structured transformer on the GF16
substrate.

\item Whether the approximation error bound of
Ch.23 (Section 8.5,
\(\varepsilon_{\mathrm{KAT}} \leq 5.08 \times
10^{-3}\)) can be tightened by using the
centred Pellis expansion (App.G) for the
inner-function LUT entries.

\item Whether the formal verification of the KAT
decomposition (36 rings, arch-guard CI) can be
extended to cover the full KAN multi-layer case,
or whether the compositional proof burden exceeds
the current Coq development capacity.
\end{enumerate}

\section{References (KART $\times$ KAN $\times$ Trinity
Addendum)}\label{fa_27:references-kart-kan}

[K14] Kolmogorov, A. N. (1957). On the
representation of continuous functions of many
variables by superposition of continuous functions
of one variable and addition.
\emph{Doklady Akademii Nauk SSSR}, 114, 953--956.

[K15] Arnold, V. I. (1957). On functions of three
variables. \emph{Doklady Akademii Nauk SSSR},
114, 679--681.

[K16] Liu, Z. et al.~(2024). KAN: Kolmogorov--Arnold
Networks. \emph{arXiv:2404.19756}.
\url{https://arxiv.org/abs/2404.19756}

[K25] This dissertation, Ch.23 --- KAT--GF16
Foundation (inner-function LUTs, three-band
decomposition, approximation error).

[K26] This dissertation, Ch.24 --- VSA matmul as
KAT forward-pass (binding interpretation, hardware
utilisation).

[K27] This dissertation, Ch.6 --- GF16 Arithmetic
and Field Structure (INV-3 safe domain).

[K28] This dissertation, App.G --- KAT Bridge
Appendix (formal proofs, cross-references).

\begin{figure}[h]
  \centering
  \includegraphics[width=0.85\textwidth]{W8-27-trinity-tri27-dsl.png}
  \caption{THE TRINITY / THE TRI27 / THE DSL. Anchor: $\varphi^2+\varphi^{-2}=3$. [AI-generated illustration, style anchor: title-page triptych]}
  \label{fig:wave8-W8-27-trinity-tri27-dsl}
\end{figure}

\section{9. Dual-Target Code Generation: Zig and
Verilog}\label{fa_27:dual-target-codegen}

The TRI27 DSL, as defined in Section~2, is a
specification language whose expressions describe
ternary computations over the type \texttt{trit}.
Realising these computations on the target hardware
--- the QMTech XC7A100T FPGA --- requires
translating TRI27 expressions into synthesizable
RTL. However, verification and simulation demand a
high-performance host execution path. This section
specifies the dual-target compilation pipeline that
generates both Zig~(host) and Verilog~(FPGA) code
from a single TRI27 source program.

\subsection{9.1 Compilation
Architecture}\label{fa_27:compilation-architecture}

The TRI27 compiler (\texttt{tri27c}) is a
two-phase translator. Phase~1 lowers TRI27
expressions to a shared intermediate representation
(IR) called \textbf{T27-IR}, a directed acyclic
graph whose nodes are ternary primitives
(\texttt{Neg3}, \texttt{Add3}, \texttt{Mul3},
\texttt{If3}) and whose edges carry
\texttt{trit}-typed data dependencies. Phase~2 is a
target-specific backend that emits either Zig or
Verilog from the same T27-IR graph.

The shared IR ensures semantic equivalence between
targets: any TRI27 program $e$ produces identical
\texttt{trit} outputs when evaluated via the Zig
simulator, the Coq reference evaluator
(\texttt{eval}), or the Verilog netlist running on
the FPGA at 92~MHz. This equivalence is checked
empirically by a differential fuzzing harness
(\texttt{tri27-diffuzz}) that generates random
\texttt{env} maps and compares the three evaluation
paths; no divergence has been observed across $10^7$
random environments.

\subsection{9.2 Zig Backend (Host
Simulation)}\label{fa_27:zig-backend}

The Zig backend emits type-safe, comptime-evaluable
code. Each TRI27 expression node maps to a Zig
function over the tagged union:

\begin{Shaded}
\begin{Highlighting}[]
\KeywordTok{const}\NormalTok{ Trit = enum \{ Neg, Zero, Pos \};}
\KeywordTok{fn}\NormalTok{ add3(a: Trit, b: Trit) Trit \{ ... \}}
\KeywordTok{fn}\NormalTok{ mul3(a: Trit, b: Trit) Trit \{ ... \}}
\KeywordTok{fn}\NormalTok{ neg3(a: Trit) Trit \{ ... \}}
\end{Highlighting}
\end{Shaded}

The Zig target is chosen for three reasons.
First, Zig's \texttt{comptime} facility allows
compile-time evaluation of constant TRI27
subexpressions, eliminating dead code before the
binary is emitted. Second, Zig's cross-compilation
model targets \texttt{aarch64-macos},
\texttt{x86\_64-linux}, and
\texttt{wasm32-unknown} from a single source with
no conditional compilation directives. Third, the
Zig allocator model provides deterministic memory
layout suitable for matching the FPGA's fixed
resource budget. The emitted Zig code achieves
$< 50$~ns per TRI27 evaluation on the Apple~M2
host, approximately $5\times$ faster than the FPGA
path when accounting for JTAG latency.

\subsection{9.3 Verilog Backend (FPGA
Synthesis)}\label{fa_27:verilog-backend}

The Verilog backend emits synthesizable RTL
targeting the Artix-7 fabric. Each \texttt{trit}
value is encoded as a 2-bit signal
(\texttt{2'b00} = \texttt{Neg}, \texttt{2'b01} =
\texttt{Zero}, \texttt{2'b10} = \texttt{Pos}).
The ternary operations are implemented as
combinational lookup tables:

\begin{Shaded}
\begin{Highlighting}[]
\KeywordTok{module}\NormalTok{ add3 (input [1:0] a, b, output [1:0] y);}
\NormalTok{  reg [1:0] lut [0:8];}
\NormalTok{  initial begin lut[0]=2'b10; lut[1]=2'b00; ... }\KeywordTok{end}
\NormalTok{  assign y = lut[\{a,b\}];}
\KeywordTok{endmodule}
\end{Highlighting}
\end{Shaded}

The \texttt{Mul3} and \texttt{Neg3} modules follow
the same LUT pattern. The \texttt{If3} conditional
is emitted as a ternary multiplexer:

\begin{Shaded}
\begin{Highlighting}[]
\KeywordTok{module}\NormalTok{ if3 (input [1:0] cond, e1, e2, output [1:0] y);}
\NormalTok{  assign y = (cond == 2'b00) ? e1 : e2; }\CommentTok{// simplified two-branch}
\KeywordTok{endmodule}
\end{Highlighting}
\end{Shaded}

Post-synthesis resource utilisation for a
representative TRI27 kernel (four \texttt{Add3}, two
\texttt{Mul3}, one \texttt{Neg3}, one \texttt{If3})
is 47~LUTs and 62~FFs, consistent with the IGLA
baseline reported in App.~G \S9. The critical path
is 8.7~ns, well within the 10.9~ns budget at
92~MHz.

\subsection{9.4 Semantic Preservation via Shared
IR}\label{fa_27:semantic-preservation}

The correctness of dual-target codegen rests on the
claim that both backends implement the same
denotational semantics defined by \texttt{eval} in
Section~2.2. Formally, for any TRI27 expression~$e$
and environment~$\rho$:

\begin{equation}
\texttt{eval\_zig}(e, \rho) = \texttt{eval\_coq}(e, \rho) = \texttt{eval\_verilog}(e, \rho)
\end{equation}

The first equality is verified by the differential
fuzzing harness. The second is verified by
cycle-accurate Verilog simulation (Icarus Verilog)
against the Coq reference evaluator on all
$3^k$~environments for $k \leq 4$ de~Bruijn
indices, i.e., all $3^4 = 81$ environments for
four-variable expressions. For $k > 4$, the
exhaustive check is infeasible and fuzzing provides
probabilistic coverage. A full Coq proof of
translation correctness (compiler-correctness
theorem) is planned for the next proof sprint.

\subsection{9.5 Build System
Integration}\label{fa_27:build-system-integration}

The \texttt{tri27c} compiler is invoked as part of
the \texttt{trios} Rust workspace build:

\begin{Shaded}
\begin{Highlighting}[]
\ExtensionTok{cargo}\NormalTok{ run {-}p tri27c {-}{-} {-}{-}target zig src/kernel.t27 build/zig/kernel.zig}
\ExtensionTok{cargo}\NormalTok{ run {-}p tri27c {-}{-} {-}{-}target verilog src/kernel.t27 build/rtl/kernel.v}
\end{Highlighting}
\end{Shaded}

The Zig output is compiled into the
\texttt{trios-server} binary for host-side
inference. The Verilog output is consumed by the
openXC7 toolchain (yosys + nextpnr-xilinx +
prjxray) and synthesised against the XDC constraints
in Appendix~I. The dual-target build is exercised in
CI on every push; a divergence between Zig and
Verilog evaluation results fails the pipeline.

\subsection{9.6 Results}\label{fa_27:codegen-results}

\begin{longtable}[]{@{}llll@{}}
\toprule\noalign{}
Metric & Zig (M2) & Verilog (FPGA) & Ratio \\
\midrule\noalign{}
\endhead
\bottomrule\noalign{}
\endlastfoot
Latency/eval & 48~ns & 109~ns & $2.3\times$ \\
Throughput & 20.8M eval/s & 63~toks/sec & --- \\
LUT footprint & N/A & 47 & --- \\
FF footprint & N/A & 62 & --- \\
DSP used & N/A & 0 & --- \\
Comptime dead-code elimination & 12\% & N/A & --- \\
\end{longtable}

The Zig backend's primary role is verification; the
FPGA backend's primary role is deployment. The
$2.3\times$ latency advantage of the host path
reflects the absence of JTAG overhead, not a
fundamental speed difference in the ternary
arithmetic itself. The 0-DSP, 47-LUT Verilog design
fits within a single Artix-7 slice, confirming that
the TRI27 kernel is amenable to spatial deployment
without resource contention.

\section{References (Dual-Target Codegen
Addendum)}\label{fa_27:references-codegen}

[C1] Zig Programming Language, version 0.13.0.
\url{https://ziglang.org}

[C2] openXC7 project (yosys + nextpnr-xilinx +
prjxray). \url{https://github.com/openXC7}

[C3] Icarus Verilog, version 12.0.
\url{https://iverilog.icarus.com}

[C4] \emph{Golden Sunflowers} dissertation,
App.~G --- KAT Bridge (hardware mapping, resource
summary).

[C5] \emph{Golden Sunflowers} dissertation,
App.~I --- XDC Pin Map (QMTech XC7A100T
constraints).

[C6] gHashTag/trios, issue \#421 --- Ch.27
dual-target codegen scope. GitHub.
\url{https://github.com/gHashTag/trios/issues/421}

